\documentclass[11pt,a4paper]{article}

% === Packages ===
\usepackage[utf8]{inputenc}
\usepackage[T1]{fontenc}
\usepackage{lmodern}
\usepackage[margin=1in]{geometry}
\usepackage{amsmath,amssymb,amsthm}
\usepackage{booktabs}
\usepackage{url}
\usepackage[numbers,sort&compress]{natbib}
\usepackage{microtype}
\usepackage{enumitem}
\usepackage{algorithm}
\usepackage{algpseudocode}
\usepackage[most]{tcolorbox}
\usepackage[hidelinks]{hyperref}

% === Errata to the published protocol, set off from the running text ===
\newtcolorbox{departure}[1]{
  breakable,
  enhanced,
  colback=black!4,
  colframe=black!60,
  boxrule=0.6pt,
  arc=1pt,
  left=6pt, right=6pt, top=4pt, bottom=4pt,
  fonttitle=\bfseries\small,
  coltitle=black,
  colbacktitle=black!10,
  title={Erratum: #1},
  before skip=8pt, after skip=8pt
}

% === Settings ===
\bibliographystyle{plainnat}
\algnewcommand{\On}[1]{\item[\textbf{on}] #1}
\algnewcommand{\Send}[2]{\State send #1 to #2}
\algnewcommand{\Guard}[1]{\State \textbf{require} #1}

% === Theorem environments, numbered by section ===
\theoremstyle{plain}
\newtheorem{theorem}{Theorem}[section]
\newtheorem{lemma}[theorem]{Lemma}
\newtheorem{corollary}[theorem]{Corollary}
\newtheorem{assumption}[theorem]{Assumption}
\theoremstyle{definition}
\newtheorem{definition}[theorem]{Definition}
\theoremstyle{remark}
\newtheorem{remark}[theorem]{Remark}

% === Interface lemmas: labelled parts ===
\newlist{iface}{description}{1}
\setlist[iface]{nosep,leftmargin=1.2em,font=\normalfont\scshape}

% === Macros for the protocol ===
\newcommand{\msg}[1]{\textsc{#1}}
\newcommand{\Inv}{\mathcal{I}}
\newcommand{\Frag}{\mathrm{Frag}}
\newcommand{\Holds}{\mathrm{Holds}}
\newcommand{\Acked}{\mathrm{Acked}}
\newcommand{\QuorumAcked}{\mathrm{QuorumAcked}}
\newcommand{\Committed}{\mathrm{Committed}}
\newcommand{\Backed}{\mathrm{Backed}}
\newcommand{\Sent}{\mathrm{Sent}}
\newcommand{\primary}{\mathrm{primary}}
\newcommand{\sent}{\mathit{sent}}
\newcommand{\started}{\mathit{started}}
\newcommand{\status}{\mathit{status}}
\newcommand{\normal}{\textsc{normal}}
\newcommand{\viewchange}{\textsc{view-change}}
\newcommand{\transfer}{\textsc{transfer}}
\newcommand{\recovering}{\textsc{recovering}}
\newcommand{\Local}{\textsc{Replica-consistency}}
\newcommand{\WellFormed}{\textsc{Message-consistency}}
\newcommand{\OneLog}{\textsc{View-agreement}}
\newcommand{\CommitsBacked}{\textsc{Commit-quorums}}
\newcommand{\Survival}{\textsc{Commit-survival}}
\newcommand{\dep}[1]{\Comment{App.~\ref{#1}}}
\newcommand{\why}[1]{\Comment{Lemma~\ref{#1}}}
\newcommand{\whydep}[2]{\Comment{Lemma~\ref{#1}; App.~\ref{#2}}}

% === Metadata ===
\title{Viewstamped Replication Made Live}
\author{Pekka Enberg}
\date{Draft of \today}

\begin{document}

\maketitle

\begin{abstract}
\noindent
Viewstamped Replication is a protocol for replicating a state machine
across a cluster of machines that can crash: a primary orders the
operations, a quorum of backups acknowledges them, and a view change
replaces the primary when it fails. It was introduced by Oki and Liskov
in 1988 and revised by Liskov and Cowling in 2012, which adds a recovery
protocol that needs no disk. The 2012 presentation has known bugs: its recovery can lose
committed data, its state transfer can lose committed data, and its
commit-number handling can execute an operation twice. This paper
formalises the protocol with the errata applied, as a transition system
and as pseudocode, and proves it safe: in every reachable state, two
replicas that have both committed an op-number hold the same entry
there, and every committed entry is held by enough replicas to be in
every quorum they could form. The proof is an inductive invariant over
the history of every message ever sent, and it is machine-checked in
Lean~4.
\\[6pt]
\textit{Keywords}: Viewstamped Replication, state machine replication,
consensus, inductive invariants, Lean
\end{abstract}

\section{Introduction}
\label{sec:intro}

Viewstamped Replication keeps a service available by running the same
deterministic operations, in the same order, on several machines. One
replica, the primary, assigns each client operation a position in a log
and sends it to the others; once a quorum has acknowledged it, the
operation is committed and every replica executes it. If the primary
fails, the backups move to a new \emph{view} with a new primary chosen
by a fixed rule, and the new primary starts from a log that a quorum has
seen, so nothing committed is lost. Oki and Liskov introduced the
protocol in 1988~\citep{oki1988vr}, inside a transactional system. Liskov
and Cowling revised it in 2012~\citep{liskov2012vr}, presenting it on its
own, with a state transfer protocol for replicas that fall behind and a
recovery protocol that lets a replica rejoin after losing all its memory
without ever having written to disk.

The 2012 presentation has known bugs. Its diskless recovery lets a
replica forget it took part in a view change and complete that change
later from stale messages, which loses committed
data~\citep{michael2017recovery}. Its state transfer truncates a log
below an entry the replica has already acknowledged, which also loses
committed data, and its commit-number handling lets a replica adopt a
lower commit number and execute an operation
twice~\citep{vanlightly2022vr}. Each has a known fix. An implementation
must apply all of them, and then runs a protocol that no single paper
states.

This paper states that protocol and proves it safe. Section~\ref{sec:prelim}
formalises the cluster as a transition system; Sections~\ref{sec:normal}
to~\ref{sec:recovery} give the handlers as pseudocode with every erratum
applied and marked; Appendix~\ref{app:vr} walks through the published
protocol and explains each correction where the paper's text leads
wrong. The proof is an inductive invariant over the history of every
message ever sent, and it is machine-checked in Lean~4
(Appendix~\ref{app:lean}).

\subsection{Our result}

Fix a cluster of $N \ge 2$ replicas running the protocol of
Sections~\ref{sec:normal} to~\ref{sec:recovery} on an asynchronous
network that may lose, delay, duplicate, and reorder messages, in which
replicas may crash and restart with nothing but their view number. Let
$k_r$ be the commit number and $L_r$ the log of replica $r$, and let
$\mathcal{R}$ be the set of replicas not currently recovering
(Section~\ref{sec:prelim}).

\begin{theorem}[Log safety]
\label{thm:log}
In every reachable cluster state:
\begin{enumerate}[nosep,label=(\roman*)]
  \item \emph{Prefix agreement.} If $n \le k_a$ and $n \le k_b$ then
    $L_a[n] = L_b[n]$.
  \item \emph{Durability.} For every $r \in \mathcal{R}$ and
    $n \le k_r$, at least $|\mathcal{R}| + 1 - Q$ replicas in
    $\mathcal{R}$ hold $L_r[n]$ at op-number $n$.
  \item \emph{Commit bounded.} $k_r \le |L_r|$.
\end{enumerate}
\end{theorem}

Durability says that every quorum the non-recovering replicas could form
contains a replica holding the committed entry; with nobody recovering
the bound is a majority. A recovering replica holds nothing and votes
for nothing, so it counts on neither side.

The proof in the body of this paper is a hand proof at the level of
detail of an algorithms paper. Every definition, lemma, and theorem in it
has a counterpart in a Lean~4 development that proves
Theorem~\ref{thm:log} in full from the same invariant, with no unproved
step; Appendix~\ref{app:lean} states the mechanised theorem, its
assumptions, and the correspondence.

\subsection{Technique}

The theorem is not preserved by a single protocol step on its own. When
an acknowledgement completes a quorum and an entry becomes committed,
whether it is now held by a majority depends on what the other
acknowledgements said, and nothing in the theorem records that. We
strengthen the statement into an inductive invariant $\Inv$
(Section~\ref{sec:invariant}) that holds initially, is preserved by
every step, and implies the theorem.

$\Inv$ is stated over the \emph{history}: the set of every message ever
sent, which delivery never removes from. Replicas forget, since a view
change replaces a backup's log, but the acknowledgement the backup sent
stays in the history, and the invariant quantifies over that. The
history also makes the network model one rule: delivery picks any
message ever sent, which covers loss, delay, duplication, and reordering
at once.

$\Inv$ has five parts, named for what they say. \Local: each replica's
own fields are consistent with one another. \WellFormed: each message's
fields agree with one another. \OneLog: within a view, an op-number
determines its entry, so every fragment of a view's log that the history
holds agrees with every other. \CommitsBacked: every commit number, on a
replica or in a message, has a quorum of acknowledgements behind every
op-number up to it. \Survival: whatever was committed in a view is held,
at the same op-number, by every later view. Around these are the clauses
the induction turned out to need, most of them about the order in which
things happened: a vote is cast after the acknowledgements its sender
made, a view is started at most once, a primary's log in its view never
shrinks. The lemma that carries the induction is the following.

\begin{lemma}[The chosen log holds earlier commits]
\label{lem:chosen-intro}
Suppose entry $e$ was acknowledged at op-number $n$ in view $v$ by a
quorum, and view $u > v$ was started from a quorum of votes $V$, and
every vote in $V$ from a view strictly between $v$ and $u$ holds $e$ at
$n$. Then the vote of $V$ from which $u$'s log was chosen holds $e$ at
$n$.
\end{lemma}

The two quorums intersect in a replica $x$. The vote $x$ cast for $u$
was cast after $x$ acknowledged $n$ in $v$, so it is from a view at least
$v$; if from $v$ itself, it covers the acknowledgement; if from a later
view, the hypothesis applies. The best vote is at least as good as
$x$'s, by the ordering the new primary uses, and holds $e$ too. A strong
induction on $u$ discharges the hypothesis. Each of the three safety
corrections in Appendix~\ref{app:vr} is what makes one step of this
argument go through: without a persisted view number a replica's view
can go backwards; without the ban on truncation a vote need not cover
what its sender acknowledged; without monotone commits \CommitsBacked{}
fails at the replica that adopted a lower number.

\subsection{Organisation}

Section~\ref{sec:prelim} fixes notation, the model of a cluster, and the
invariant. Section~\ref{sec:framework} states the interface of every
handler as a lemma and proves the theorem from those interfaces, before
any handler is shown. Sections~\ref{sec:normal} to~\ref{sec:recovery}
give the handlers as pseudocode, one section per sub-protocol, each with
an overview and the analysis that proves its interface lemmas.
Appendix~\ref{app:vr} is the published protocol and its corrections.
Appendix~\ref{app:lean} is the mechanised proof.

\section{Preliminaries}
\label{sec:prelim}

\subsection{The cluster}

\begin{definition}[Configuration]
\label{def:config}
$N \ge 2$ replicas with identities $0, \dots, N-1$; $f = \lfloor N/2
\rfloor$; the quorum size $Q = f + 1$; and $\primary(v) = v \bmod N$, the
primary of view $v$.
\end{definition}

\begin{definition}[Log]
A log $L$ is a finite sequence of entries $e = (c, s, \mathit{op})$,
client identity, request number, operation. The position of an entry in
the log is its \emph{op-number}, counted from $1$: $L[n]$ is the entry
with op-number $n$, for $1 \le n \le |L|$. $L(a{..}b]$ is the entries
with op-numbers $a+1, \dots, b$, and $L \cdot e$ is the log with $e$
appended, at op-number $|L|+1$.
\end{definition}

\begin{definition}[Replica]
\label{def:replica}
A replica $r$ has an identity $\mathrm{id}_r < N$; a status
$\status_r$, one of \normal, \viewchange, \transfer, or \recovering; a view
number $v_r$; a last normal view $\ell_r$, the latest view in which
$\status_r$ was \normal; a log $L_r$ with $n_r = |L_r|$; a commit number
$k_r$; a flag $\mathit{catching}_r$; a client table $C_r$ mapping client
identities to a request number and an optional reply; acknowledgements
$A_r$ mapping op-numbers to sets of identities; a set $S_r$ of
identities; a flag $\mathit{voted}_r$; votes $V_r$ mapping identities to
votes; a nonce $x_r$; recovery responses $R_r$ mapping identities to
responses; and a state machine state $\sigma_r$.
$\mathrm{isPrimary}(r)$ means $\mathrm{id}_r = \primary(v_r)$.
\end{definition}

\begin{definition}[Messages]
\label{def:messages}
Eleven kinds. $i$ is an identity, $v$ a view, $n$ an op-number, $k$ a
commit number, $L$ a log, $x$ a nonce, and $\mathit{st} \in \{\bot\} \cup
\{(L, k)\}$ an optional state.
\begin{center}
\small
\begin{tabular}{@{}ll@{}}
$\msg{Request}(c, s, \mathit{op})$ & $\msg{StartViewChange}(v, i)$ \\
$\msg{Prepare}(v, n, e, k)$ & $\msg{DoViewChange}(v, i, \ell, L, n, k)$ \\
$\msg{PrepareOk}(v, n, i)$ & $\msg{StartView}(v, L, n, k)$ \\
$\msg{Commit}(v, k)$ & $\msg{Recovery}(i, x, v)$ \\
$\msg{GetState}(i, v, n)$ & $\msg{RecoveryResponse}(v, x, i, \mathit{st})$ \\
$\msg{NewState}(v, L, a, b, k)$ &
\end{tabular}
\end{center}
A \emph{vote} is the triple $(\ell, L, k)$ a \msg{DoViewChange} carries:
the sender's last normal view, its log, its commit number.
\end{definition}

\begin{definition}[Cluster]
\label{def:cluster}
A cluster state is $\Sigma = (\mathit{replicas}, \sent, \mathit{replies},
\started)$: the replicas by identity; $\sent$, the set of every
$(\mathit{to}, m)$ ever sent; $\mathit{replies}$, every reply
$(v, c, s, x)$ ever produced; and $\started$, the set of $(v, V)$ for
every view $v$ some primary started and the votes $V$ it chose the log
from. $\Sent(\mathit{to}, m)$ means $(\mathit{to}, m) \in \sent$.
$\mathcal{R}$ is the set of replicas with $\status \ne \recovering$.
\end{definition}

\begin{definition}[Step]
\label{def:step}
$\Sigma \to \Sigma'$ is one of: \emph{deliver} some $(\mathit{to}, m) \in
\sent$ to replica $\mathit{to}$; \emph{idle} at some replica; \emph{request}
$(c, s, \mathit{op})$ at some replica; \emph{recover} some replica with
the view number it holds in $\Sigma$ and a nonce that no \msg{Recovery}
in $\sent$ carries. The replica runs the handler of
Sections~\ref{sec:normal} to~\ref{sec:recovery}; everything it sends,
replies, and starts is added to $\sent$, $\mathit{replies}$, and
$\started$, which therefore only grow.
\end{definition}

\begin{definition}[Reachable]
\label{def:reachable}
The initial cluster has every replica with $\status = \normal$,
$v = \ell = k = 0$, $L$ empty, every map and set empty, $\sigma$ the
initial state, and the three sets of $\Sigma$ empty. A cluster is
\emph{reachable} if it is obtained from the initial cluster by finitely
many steps.
\end{definition}

\begin{assumption}
\label{asm:model}
Throughout:
\begin{enumerate}[nosep,label=(\roman*)]
  \item $N \ge 2$.
  \item The environment chooses which replica takes an idle step and
    when; the protocol's timers count those idle steps. Every schedule is
    a reachable execution.
  \item A replica's view number $v_r$ is written to stable storage after
    every step and before anything the step produced is sent, and is the
    only state a \emph{recover} step retains.
  \item The nonce a \emph{recover} step uses is fresh: no \msg{Recovery}
    in the history carries it.
\end{enumerate}
\end{assumption}

\begin{remark}
(i) excludes only the one-replica cluster, which never sends a message,
so nothing about it can be stated through $\sent$; it is trivially safe.
(ii) means safety is proved for every timing, including ones no timer
would produce. (iii) is the one erratum to the published protocol's
diskless design, and Appendix~\ref{app:recovery} shows it cannot be
dropped. (iv) is the paper's ``nonce'' (\S4.3), which a clock or a
persisted counter provides.
\end{remark}

\begin{lemma}[Quorum intersection]
\label{lem:quorum}
Two sets of at least $Q$ identities below $N$ have a common element.
More precisely, two sets $X, Y$ of identities below $N$ share at least
$|X| + |Y| - N$ elements.
\end{lemma}
\begin{proof}
$|X \cap Y| = |X| + |Y| - |X \cup Y| \ge |X| + |Y| - N$, and
$2Q = 2\lfloor N/2 \rfloor + 2 > N$.
\end{proof}

Delivery chooses a message from $\sent$ and leaves it there, so a
message may arrive any number of times, in any order, or never. A
replica that is never chosen to run is a crashed one.

\subsection{The history}

The invariant is stated over the history. The following definitions name
the pieces of a view's log that the history holds.

\begin{definition}[Vote of]
\label{def:voteof}
$\mathrm{Vote}(u, x, d)$: replica $x$ cast the vote $d = (\ell, L, k)$
for view $u$. Either the history holds
$\msg{DoViewChange}(u, x, \ell, L, |L|, k)$, or $(u, V) \in \started$
with $(x, d) \in V$. The second case covers the new primary's own vote,
which it records without sending.
\end{definition}

\begin{definition}[Fragment]
\label{def:frag}
$\Frag(v, a, L)$: the history holds $L$ as a piece of view $v$'s log,
the entries with op-numbers $a+1, \dots, a+|L|$. One rule per kind of
message that carries log:
\[
\begin{array}{ll}
\Sent(\_, \msg{Prepare}(v, n, e, \_)) & \Rightarrow \Frag(v, n-1, [e]) \\
\Sent(\_, \msg{NewState}(v, L, a, \_, \_)) & \Rightarrow \Frag(v, a, L) \\
\Sent(\_, \msg{StartView}(v, L, \_, \_)) & \Rightarrow \Frag(v, 0, L) \\
\Sent(\_, \msg{RecoveryResponse}(v, \_, \_, (L, \_))) & \Rightarrow \Frag(v, 0, L) \\
\mathrm{Vote}(\_, \_, (\ell, L, \_)) & \Rightarrow \Frag(\ell, 0, L)
\end{array}
\]
\end{definition}

\begin{remark}
A vote is a fragment of the view its sender was \emph{last normal in},
$\ell$, not of the view it votes for. A replica catching up with a view
has a view number ahead of its log; tagging by $\ell$ is what keeps it
consistent with the rest.
\end{remark}

\begin{definition}[Holds]
$\Holds(v, n, e)$: there are $a < n$ and $L$ with $\Frag(v, a, L)$ and
$L[n - a] = e$.
\end{definition}

\begin{definition}[Acknowledged]
\leavevmode
\begin{enumerate}[nosep,label=(\roman*)]
  \item $\Acked(v, n, q)$: $q = \primary(v)$, or the history holds
    $\msg{PrepareOk}(v, o, q)$ for some $o \ge n$.
  \item $\QuorumAcked(v, n)$: there are $Q$ distinct identities $q < N$
    with $\Acked(v, n, q)$.
\end{enumerate}
\end{definition}

\begin{definition}[Committed]
$\Committed(v, n, e)$: $\Holds(v, n, e)$ and $\QuorumAcked(v, n)$.
\end{definition}

\begin{definition}[Backed]
$\Backed(v, k)$: for every op-number $n \le k$,
\[
\exists\, e,\ v' \le v.\quad \Committed(v', n, e) \wedge \Holds(v, n, e).
\]
\end{definition}

\begin{lemma}[The history is monotone]
\label{lem:mono}
$\mathrm{Vote}$, $\Frag$, $\Holds$, $\Acked$, $\Committed$, and $\Backed$
are preserved by every step.
\end{lemma}
\begin{proof}
Each is defined by the existence of elements of $\sent$ and $\started$,
which only grow.
\end{proof}

\subsection{The invariant}
\label{sec:invariant}

\begin{definition}[\Local]
\label{def:local}
At every replica: $k_r \le n_r$; $\ell_r \le v_r$; if $\status_r$ is
\normal{} or \transfer{} then $\ell_r = v_r$; if $\mathit{catching}_r$
then $\status_r = \viewchange$; if $\status_r = \recovering$ then
$L_r = \varepsilon$ and $k_r = 0$.
\end{definition}

\begin{definition}[\WellFormed]
\label{def:wellformed}
Every message in $\sent$ has fields that agree with one another: a
\msg{StartView} or \msg{DoViewChange} log is as long as its op-number and
its commit number does not exceed it; a \msg{NewState} log is as long as
$b - a$ with $a \le b$ and $k \le b$; a \msg{DoViewChange} has
$\ell \le v$; a \msg{Prepare} has $n > 0$; a \msg{RecoveryResponse} state
has $k \le |L|$.
\end{definition}

\begin{definition}[\OneLog]
\label{def:viewagreement}
\leavevmode
\begin{enumerate}[nosep,label=(\roman*)]
  \item Within a view, an op-number determines its entry:
    $\Holds(v, n, e) \wedge \Holds(v, n, e') \Rightarrow e = e'$.
  \item Every replica's log agrees with the fragments of the view it
    came from: $L_r[n] = e \wedge \Holds(\ell_r, n, e') \Rightarrow e = e'$.
\end{enumerate}
\end{definition}

\begin{remark}
\OneLog{} says nothing across views: op-number $n$ may name different
entries in views $4$ and $5$, which is the paper's own observation
(\S4.2) that two operations can be assigned the same op-number only
across a view change, and \Survival{} governs which of them matter. It
is the property Raft calls log matching, with views for terms and message
fragments alongside replica logs, and it is what lets the rest of the
invariant say ``the entry at op-number $n$ in view $v$'' and mean one
thing.
\end{remark}

\begin{definition}[\CommitsBacked]
$\Backed(\ell_r, k_r)$ for every replica; and for every message in
$\sent$ carrying a commit number $k$ in view $v$, $\Backed(v, k)$, a vote
being judged by its last normal view.
\end{definition}

\begin{definition}[\Survival]
If $\Committed(v', n, e)$ and $v' < v$, then every \msg{StartView} log
of $v$, every vote from $v$, every recovery state of $v$, every
\msg{NewState} segment of $v$ that reaches $n$, every \msg{Prepare} of
$v$ for op-number $n$, and every non-recovering replica with
$\ell_r = v$ holds $e$ at $n$.
\end{definition}

\begin{remark}
This is the correctness condition of \citet[\S8.1]{liskov2012vr}, that
``every committed operation survives into all subsequent views in the
same position in the serial order'', stated over every fragment of a
view rather than over the view as a whole.
\end{remark}

\begin{definition}[The invariant $\Inv$]
\label{def:inv}
\Local, \WellFormed, \OneLog, \CommitsBacked, and \Survival, together
with the following clauses. Appendix~\ref{app:lean} lists each by its
name in the mechanised proof.
\begin{enumerate}[nosep,label=(\alph*)]
  \item \emph{Views only rise.} No message in $\sent$, and no vote,
    carries a view above its sender's current view. A replica in
    \viewchange{} status has $\ell_r < v_r$; a vote for $u$ is from a
    view $\ell < u$; the primary of a started view $v$ has $\ell \ge v$.
  \item \emph{Chosen from a quorum, once.} A view is started at most
    once. A started view's \msg{StartView} log extends the best of the
    $Q$ votes, from distinct replicas, it was started from. Every whole
    log of a started view, and every \msg{NewState} of it, reaches at
    least as far as that best vote's log.
  \item \emph{Votes cover acknowledgements.} A vote from $x$ covers
    every op-number $x$ acknowledged in the vote's view $\ell$; if $x$ is
    the primary of $\ell$, it covers every fragment of $\ell$.
  \item \emph{Votes come after acknowledgements.} If $x$ acknowledged in
    view $v$, or is the primary of a started view $v$, and cast a vote
    for $u > v$, that vote is from a view $\ell \ge v$.
  \item \emph{Primary is longest.} While $v$ is the last normal view of
    $\primary(v)$ and it is not recovering, every fragment of $v$, every
    log of a replica last normal in $v$, and every acknowledgement in
    $v$ is within its log.
  \item \emph{Acknowledgements are kept.} If $q$ sent
    $\msg{PrepareOk}(v, o, q)$ then $\ell_q \ge v$, and while
    $\ell_q = v$ and $q$ is not recovering, $n_q \ge o$. A normal
    primary's record $A_r$ holds only acknowledgements of its own
    op-numbers in its own view, each acknowledger once.
  \item \emph{Recovery covers acknowledgements.} A recovery state is the
    log of the primary of its view, answers a \msg{Recovery} with its
    nonce, and, when it answers a replica now recovering, covers every
    op-number that replica acknowledged in the state's view.
  \item \emph{Bookkeeping.} Replica $i$ of the cluster has identity $i$;
    there are $N$ replicas; every message names a sender below $N$; a
    replica catching up or in state transfer is not its view's primary;
    no \msg{Prepare} or \msg{Commit} is addressed to its view's primary;
    every entry of every log is held by some fragment of the replica's
    last normal view; two replicas with the same last normal view agree;
    every view above zero with a fragment, an acknowledgement, or a
    non-recovering replica was started, and every started view has a
    \msg{StartView}; what a replica has recorded in $V_r$ or $R_r$ was
    sent to it; between steps every outbox is empty; and no internal
    guard of \Call{InstallLog}{} or \Call{CommitUpTo}{} has failed.
\end{enumerate}
\end{definition}

\begin{remark}
The three safety corrections of Appendix~\ref{app:vr} are each one
clause. Truncation in state transfer would falsify (c). Diskless
recovery would falsify (a), since a recovered replica's view could go
down. A decreasing commit number would falsify \CommitsBacked{} at the
replica that adopted it.
\end{remark}

Table~\ref{tab:notation} collects the notation.

\begin{table}[t]
\centering
\small
\caption{Notation. Subscript $r$ names the replica; it is dropped when
only one replica is in view.}
\label{tab:notation}
\begin{tabular}{@{}ll@{}}
\toprule
Symbol & Meaning \\
\midrule
$N$ & number of replicas; identities are $0, \dots, N-1$ \\
$f$ & $\lfloor N/2 \rfloor$, the failures tolerated \\
$Q$ & $f + 1$, the quorum size \\
$\primary(v)$ & $v \bmod N$, the primary of view $v$ \\
\addlinespace
$v_r$ & view number of replica $r$ \\
$\ell_r$ & last normal view: the latest view $r$ was normal in \\
$L_r$ & log of $r$, a sequence of entries with op-numbers $1, \dots, |L_r|$ \\
$n_r$ & $|L_r|$, the op-number of the latest entry \\
$k_r$ & commit number: the entries with op-numbers $1, \dots, k_r$ are committed \\
$\sigma_r$ & state of $r$'s state machine \\
$A_r, S_r, V_r, R_r, C_r$ & acknowledgements, view-change senders, votes, recovery responses, client table \\
\addlinespace
$e = (c, s, \mathit{op})$ & a log entry: client, request number, operation \\
$L[n]$, $L(a{..}b]$, $L \cdot e$ & entry with op-number $n$; op-numbers $a+1, \dots, b$; append \\
$(\ell, L, k)$ & a vote: last normal view, log, commit number \\
\addlinespace
$\Sigma$ & a cluster state \\
$\sent$ & every $(\mathit{to}, m)$ ever sent; never shrinks \\
$\started$ & every $(v, V)$: view $v$ was started from votes $V$ \\
$\mathcal{R}$ & the replicas whose status is not recovering \\
$\Sent(\mathit{to}, m)$ & $(\mathit{to}, m) \in \sent$ \\
$\mathrm{Vote}(u, x, d)$ & $x$ cast vote $d$ for view $u$, sent or recorded \\
\addlinespace
$\Frag(v, a, L)$ & the history holds $L$ as view $v$'s entries after op-number $a$ \\
$\Holds(v, n, e)$ & some fragment of view $v$ has $e$ at op-number $n$ \\
$\Acked(v, n, q)$ & replica $q$ acknowledged op-number $n$ in view $v$ \\
$\QuorumAcked(v, n)$ & $Q$ distinct replicas acknowledged $n$ in $v$ \\
$\Committed(v, n, e)$ & $\Holds(v, n, e)$ and $\QuorumAcked(v, n)$ \\
$\Backed(v, k)$ & every op-number up to $k$ is committed with the entry $v$ holds \\
$\Inv$ & the invariant, Definition~\ref{def:inv} \\
\bottomrule
\end{tabular}
\end{table}

\section{The framework}
\label{sec:framework}

This section states what each handler guarantees and proves the theorem
from those guarantees. The handlers themselves are in
Sections~\ref{sec:normal} to~\ref{sec:recovery}, each of which ends by
proving the interface lemmas stated here.

\paragraph{What is assumed and what is proved.} The line is drawn
where the protocol's control ends. Assumption~\ref{asm:model} is about
the world the protocol runs in: the network, what a crash keeps, when
timers fire, where a nonce comes from. Nothing the protocol produces is
assumed. A replica's fields are written by its handlers, and a message's
fields by the replica that sends it, so any fact about them is the
protocol's responsibility and must come out of the induction. That
includes the two facts every handler relies on to run at all: that its
own commit number is within its own log, and that the message it was
handed has fields that agree with each other. These are the input
requirements of the handlers, in the sense that a subroutine has
preconditions its callers must meet. The callers here are the handler
that ran before, for a replica's state, and the sender, for a message.
We do not assume the callers comply. We prove it at every call site:
\Local{} is proved once for every handler, and \WellFormed{} once for
every send. Everything above them in the invariant is proved the same
way, one handler at a time, from the interface lemmas below.

\paragraph{Where a commit is born.} $\Committed(v, n, e)$ is a fact
about the history: $Q$ acknowledgements of $n$ in $v$ exist. It
therefore comes into existence at the moment the $Q$th
\msg{PrepareOk} is \emph{sent}, not when the primary receives it, and
the obligation to show that every later view holds $e$
(\Survival) falls on the sender. Three handlers send acknowledgements,
\msg{Prepare}, \msg{NewState}, and \msg{StartView}, and each of their
interface lemmas carries that obligation. The primary's handler for
\msg{PrepareOk} only reads a quorum that already exists.

\subsection{Interfaces of the handlers}

Every interface lemma has the same shape. \textsc{Precondition} is what
the handler may assume, always including $\Inv$. \textsc{Effect} is what
it changes and sends. \textsc{Guarantee} is what holds afterwards:
$\Inv$ again, plus the fact the step establishes that the rest of the
proof uses. A handler whose guard fails has no effect and trivially
preserves $\Inv$.

One guarantee is shared by every handler and needs no assumption about
the message received, so it is stated once.

\begin{lemma}[Local]
\label{lem:local}
Every handler of Sections~\ref{sec:normal} to~\ref{sec:recovery}, given
any message, leaves the replica satisfying \Local{}
(Definition~\ref{def:local}): its commit number within its log, its
last normal view not ahead of its view, and its status consistent with
both. Every message it sends satisfies \WellFormed{}
(Definition~\ref{def:wellformed}).
\end{lemma}
\begin{proof}
$k_r$ rises only in \Call{CommitUpTo}{}, which stops at $n_r$; a log is
replaced only in \Call{InstallLog}{}, which refuses one shorter than
$k_r$; $\ell_r$ is set only in \Call{EnterNormal}{}, to $v_r$; $v_r$ is
set only to a larger value or, on recovery, to the persisted value.
Every log sent is sent with its own length and the sender's commit
number, which \Local{} bounds.
\end{proof}

The remaining lemmas are one per handler.

\begin{lemma}[Request]
\label{lem:h-request}
\begin{iface}
  \item[Precondition] $\Inv$; the replica is the normal primary of $v_r$
    and the request is new for its client.
  \item[Effect] Appends $e$ at op-number $n_r + 1$, records its own
    acknowledgement of it, and sends $\msg{Prepare}(v_r, n_r+1, e, k_r)$
    to the others.
  \item[Guarantee] $\Inv$. The new fragment $\Frag(v_r, n_r, [e])$ lies
    beyond every existing fragment of $v_r$.
\end{iface}
\end{lemma}

\begin{lemma}[Prepare]
\label{lem:h-prepare}
\begin{iface}
  \item[Precondition] $\Inv$; the message is accepted: $v = v_r$ and the
    replica is a normal backup.
  \item[Effect] If $n = n_r + 1$, appends $e$; commits up to
    $\min(k, n_r)$; sends $\msg{PrepareOk}(v_r, n_r, \mathrm{id}_r)$ to
    the primary. If $n > n_r + 1$, enters state transfer instead.
  \item[Guarantee] $\Inv$. In particular, if the acknowledgement
    completes a quorum for some op-number $m \le n_r$ in $v_r$, then
    $\Committed(v_r, m, L_r[m])$ now holds and every started view above
    $v_r$ holds $L_r[m]$ at $m$.
\end{iface}
\end{lemma}

\begin{lemma}[PrepareOk]
\label{lem:h-prepareok}
\begin{iface}
  \item[Precondition] $\Inv$; the replica is the normal primary of
    $v_r$, $n > k_r$, and $n \in \mathrm{dom}\, A_r$.
  \item[Effect] Adds $i$ to $A_r[n]$; if that makes $|A_r[n]| = Q$,
    commits up to $n$ and replies to the clients of the committed
    entries.
  \item[Guarantee] $\Inv$. If the quorum completed, then
    $\QuorumAcked(v_r, m)$ already holds in the history for every
    $m \le n$, and hence $\Backed(v_r, n)$.
\end{iface}
\end{lemma}

\begin{lemma}[Commit]
\label{lem:h-commit}
\begin{iface}
  \item[Precondition] $\Inv$; the message is accepted.
  \item[Effect] Commits up to $k$, or enters state transfer if
    $k > n_r$.
  \item[Guarantee] $\Inv$.
\end{iface}
\end{lemma}

\begin{lemma}[GetState]
\label{lem:h-getstate}
\begin{iface}
  \item[Precondition] $\Inv$; the replica is normal in $v = v_r$ and
    $n \le n_r$.
  \item[Effect] Sends $\msg{NewState}(v, L_r(n{..}n_r], n, n_r, k_r)$ to $i$.
  \item[Guarantee] $\Inv$. The new fragment $\Frag(v, n, L_r(n{..}n_r])$
    agrees with every fragment of $v$.
\end{iface}
\end{lemma}

\begin{lemma}[NewState]
\label{lem:h-newstate}
\begin{iface}
  \item[Precondition] $\Inv$; $v = v_r$.
  \item[Effect] In \transfer{} status, appends the part of $L$ past
    $n_r$, commits up to $k$, becomes normal, and acknowledges its whole
    log. In \viewchange{} status while catching up with $a = k_r$,
    replaces $L_r(k_r{..}n_r]$ by $L$, commits up to $k$, becomes normal
    in $v_r$, and acknowledges its whole log.
  \item[Guarantee] $\Inv$. Afterwards $\ell_r = v_r$ and $L_r$ agrees
    with every fragment of $v_r$. The acknowledgement carries the same
    guarantee as in Lemma~\ref{lem:h-prepare}.
\end{iface}
\end{lemma}

\begin{lemma}[StartViewChange and DoViewChange]
\label{lem:h-svc}
\begin{iface}
  \item[Precondition] $\Inv$; $v \ge v_r$.
  \item[Effect] If $v > v_r$, adopts $v$ in \viewchange{} status and
    announces it. Records the sender. On $f$ senders, casts its vote
    $(\ell_r, L_r, k_r)$: sends it to $\primary(v_r)$, or records it if
    it is the primary. On a \msg{DoViewChange} as primary, records the
    vote; on the $Q$th, starts the view (Lemma~\ref{lem:h-start}).
  \item[Guarantee] $\Inv$. A vote cast by $r$ satisfies clauses (a), (c),
    and (d) of Definition~\ref{def:inv}: it is from $\ell_r < v_r$, it
    covers every op-number $r$ acknowledged in $\ell_r$, and it comes
    after every acknowledgement $r$ made.
\end{iface}
\end{lemma}

\begin{lemma}[Starting a view]
\label{lem:h-start}
\begin{iface}
  \item[Precondition] $\Inv$; the replica is $\primary(v_r)$ in
    \viewchange{} status and holds $Q$ votes $V$ from distinct replicas,
    each sent to it for $v_r$ or its own.
  \item[Effect] Adds $(v_r, V)$ to $\started$; installs the log $L$ of
    the best vote; commits up to the greatest $k$ in $V$; becomes normal
    in $v_r$; sends $\msg{StartView}(v_r, L, |L|, k_r)$ to the others.
  \item[Guarantee] $\Inv$. For every $u < v_r$ and every
    $\Committed(u, n, e)$: $L[n] = e$. The guards of
    \Call{InstallLog}{} and \Call{CommitUpTo}{} hold.
\end{iface}
\end{lemma}

\begin{lemma}[StartView]
\label{lem:h-startview}
\begin{iface}
  \item[Precondition] $\Inv$; $v > v_r$, or $v = v_r$ in \viewchange{}
    status.
  \item[Effect] Adopts $v$; installs $L$; commits up to $k$; becomes
    normal in $v$; acknowledges its whole log.
  \item[Guarantee] $\Inv$. Afterwards $\ell_r = v$ and $L_r = L$. The
    acknowledgement carries the same guarantee as in
    Lemma~\ref{lem:h-prepare}. The guard of \Call{InstallLog}{} holds.
\end{iface}
\end{lemma}

\begin{lemma}[Recovery and RecoveryResponse]
\label{lem:h-recovery}
\begin{iface}
  \item[Precondition] $\Inv$.
  \item[Effect] On \msg{Recovery} carrying $v > v_r$, adopts $v$. On
    \msg{Recovery} in normal status, answers with $v_r$ and, if primary,
    with $(L_r, k_r)$. On the $Q$th \msg{RecoveryResponse} for its nonce,
    if one is from the primary of the latest view $v^*$ among them and
    $v^* \ge v_r$, adopts that primary's log and commit number and becomes
    normal in $v^*$.
  \item[Guarantee] $\Inv$. A recovery state sent is a fragment of its
    view that covers every acknowledgement in that view, clause (g).
\end{iface}
\end{lemma}

\begin{lemma}[Recover]
\label{lem:h-recover}
\begin{iface}
  \item[Precondition] $\Inv$; the persisted view is $v_r$; the nonce is
    fresh (Assumption~\ref{asm:model}(iv)).
  \item[Effect] Resets the replica to its initial state except
    $v_r = \ell_r = v_r$, in \recovering{} status with the fresh nonce,
    and sends $\msg{Recovery}(\mathrm{id}_r, x, v_r)$ to the others.
  \item[Guarantee] $\Inv$. In particular clause (a): no message the
    replica ever sent is ahead of $v_r$; and clause (g) for the new
    nonce, vacuously, since no response to it exists yet.
\end{iface}
\end{lemma}

\begin{lemma}[Idle]
\label{lem:h-idle}
\begin{iface}
  \item[Precondition] $\Inv$.
  \item[Effect] Re-sends: the primary its \msg{Commit} and every
    uncommitted \msg{Prepare}; a recovering replica its \msg{Recovery};
    a replica in transfer its \msg{GetState}; a replica in view change
    its \msg{StartViewChange} and, if it has voted, its vote. A backup or
    a replica in view change whose timer fires starts view $v_r + 1$.
  \item[Guarantee] $\Inv$. Every fragment re-sent is one the history
    already holds.
\end{iface}
\end{lemma}

\subsection{The invariant is inductive}

\begin{lemma}[Initial]
\label{lem:init}
$\Inv$ holds in the initial cluster.
\end{lemma}
\begin{proof}
Nothing has been sent, so every clause about the history is vacuous;
every replica is fresh, so every clause about replica state holds.
\end{proof}

\begin{theorem}[$\Inv$ is inductive]
\label{thm:inductive}
$\Inv$ holds in the initial cluster, and if $\Inv(\Sigma)$ and
$\Sigma \to \Sigma'$ then $\Inv(\Sigma')$. Hence $\Inv$ holds in every
reachable cluster.
\end{theorem}
\begin{proof}
Lemma~\ref{lem:init} for the initial cluster. A step runs one handler at
one replica; every handler is covered by exactly one of
Lemmas~\ref{lem:h-request} to~\ref{lem:h-idle}, whose guarantee includes
$\Inv$, or its guard fails and nothing changes. Induction on the number
of steps.
\end{proof}

\subsection{Proof of Theorem~\ref{thm:log}}

\begin{lemma}[Committed entries are what the view holds]
\label{lem:committed}
If $\Inv(\Sigma)$ and $n \le k_r$, there are $e$ and $v' \le \ell_r$ with
$L_r[n] = e$, $\Committed(v', n, e)$, and $\Holds(\ell_r, n, e)$.
\end{lemma}
\begin{proof}
\CommitsBacked{} gives $e$, $v'$, and $\Holds(\ell_r, n, e)$; \OneLog{}
identifies $e$ with $L_r[n]$.
\end{proof}

\begin{lemma}[Later views hold earlier commits]
\label{lem:later}
If $\Inv(\Sigma)$, $\Committed(v', n, e)$, $v' < v$, and
$\Holds(v, n, e')$, then $e' = e$.
\end{lemma}
\begin{proof}
Cases on the fragment that holds $e'$, each a clause of \Survival.
\end{proof}

\begin{lemma}[An acknowledger holds the entry]
\label{lem:acker}
If $\Inv(\Sigma)$, $\Committed(v, n, e)$, and $z \in \mathcal{R}$ with
$\Acked(v, n, \mathrm{id}_z)$, then $L_z[n] = e$.
\end{lemma}
\begin{proof}
First $\ell_z \ge v$: by clause (f) if $z$ sent the acknowledgement, and
by clause (a) if $z$ is the primary of $v$, since $v$ was started
(clause (h)) and its primary has $\ell \ge v$. If $\ell_z > v$, the
replica clause of \Survival{} gives $L_z[n] = e$. If $\ell_z = v$, the
log reaches $n$: by clause (f) for a sent acknowledgement, by clause
(e) for the primary, whose log covers every fragment of $v$ including the
one holding $e$; and \OneLog(ii) identifies $L_z[n]$ with $e$.
\end{proof}

\begin{proof}[Proof of Theorem~\ref{thm:log}]
Let $\Sigma$ be reachable; $\Inv(\Sigma)$ by
Theorem~\ref{thm:inductive}.

\emph{Prefix agreement.} For $n \le k_a, k_b$,
Lemma~\ref{lem:committed} gives $e_a$ committed in $v_a$ and $e_b$
committed in $v_b$; if $v_a < v_b$ then $e_b = e_a$ by
Lemma~\ref{lem:later}, symmetrically if $v_b < v_a$, and by \OneLog(i)
if equal.

\emph{Durability.} Let $r \in \mathcal{R}$ and $n \le k_r$. By
Lemma~\ref{lem:committed}, $e = L_r[n]$ is committed in some $v$, so
some set $X$ of $Q$ distinct identities acknowledged $n$ in $v$. Let $P$
be the identities of $\mathcal{R}$. By Lemma~\ref{lem:quorum},
$|X \cap P| \ge Q + |\mathcal{R}| - N \ge |\mathcal{R}| + 1 - Q$, using
$2Q \ge N + 1$. Every replica in $X \cap P$ holds $e$ at $n$ by
Lemma~\ref{lem:acker}.

\emph{Commit bounded.} \Local.
\end{proof}

\section{Normal operation and state transfer}
\label{sec:normal}

\subsection{Overview}

The primary orders requests by appending them to its log and telling the
backups with \msg{Prepare}. A backup acknowledges with its whole log
length, so one \msg{PrepareOk} covers every op-number up to it. The
primary commits an op-number, and everything before it, on $Q$
acknowledgements counting its own, and tells the backups through the
commit number in the next \msg{Prepare} or an idle \msg{Commit}. A
backup that finds a gap asks for the missing log with \msg{GetState}; a
replica that learns of a later view keeps its committed prefix and asks
for the rest of the new view's log from there.
Algorithm~\ref{alg:helpers} gives the helpers every section uses;
Algorithm~\ref{alg:normal} the handlers.

\begin{algorithm}[tp]
\caption{Log and status helpers. $r$ is the replica.}
\label{alg:helpers}
\begin{algorithmic}[1]
\Procedure{Append}{$e$}
  \State $L_r \gets L_r \cdot e$;\quad $C_r[e.c] \gets (e.s, \bot)$
\EndProcedure
\Procedure{InstallLog}{$L$}
  \Guard{$|L| \ge k_r$} \why{lem:local}
  \State $L_r \gets L$; rebuild $C_r$ from $L$, keeping replies for op-numbers up to $k_r$
\EndProcedure
\Procedure{CommitUpTo}{$k$, $\mathit{reply}$}
  \While{$k_r < k$} \dep{app:normal}
    \Guard{$k_r < n_r$} \why{lem:local}
    \State $k_r \gets k_r + 1$
    \State $(\sigma_r, x) \gets \mathrm{apply}(\sigma_r, L_r[k_r].\mathit{op})$;\quad record $x$ in $C_r[L_r[k_r].c]$
    \If{$\mathit{reply}$} reply $(v_r, L_r[k_r].c, L_r[k_r].s, x)$ \EndIf
  \EndWhile
\EndProcedure
\Procedure{EnterNormal}{}
  \State $\status_r \gets \normal$;\ $\ell_r \gets v_r$;\ $S_r, V_r \gets \emptyset$;\ $\mathit{voted}_r, \mathit{catching}_r \gets \mathrm{false}$
\EndProcedure
\Procedure{StateTransfer}{}
  \State $\status_r \gets \transfer$
  \Send{$\msg{GetState}(\mathrm{id}_r, v_r, n_r)$}{the primary}
\EndProcedure
\Procedure{CatchUp}{$v$}
  \If{$v = v_r \wedge \mathit{catching}_r$} \Return \EndIf
  \State $v_r \gets v$;\ $\status_r \gets \viewchange$;\ $\mathit{catching}_r \gets \mathrm{true}$;\ $S_r, V_r \gets \emptyset$;\ $\mathit{voted}_r \gets \mathrm{false}$
  \Send{$\msg{GetState}(\mathrm{id}_r, v, k_r)$}{$\primary(v)$} \dep{app:transfer}
\EndProcedure
\Procedure{Accept}{$v$} \Comment{may this normal-case message be processed?}
  \If{$v < v_r$} \Return false \EndIf
  \If{$v > v_r$} \Call{CatchUp}{$v$}; \Return false \EndIf
  \If{$\status_r = \normal$} \Return $\lnot\mathrm{isPrimary}(r)$ \EndIf
  \If{$\status_r = \viewchange$} \Call{CatchUp}{$v$} \EndIf
  \State \Return false
\EndProcedure
\end{algorithmic}
\end{algorithm}

\begin{algorithm}[tp]
\caption{Normal operation and state transfer.}
\label{alg:normal}
\begin{algorithmic}[1]
\On{$\msg{Request}(c, s, \mathit{op})$} \why{lem:h-request}
  \Guard{$\mathrm{isPrimary}(r) \wedge \status_r = \normal$}
  \If{$C_r[c]$ undefined or $s > C_r[c].s$}
    \State \Call{Append}{$(c, s, \mathit{op})$};\ $A_r[n_r] \gets \{\mathrm{id}_r\}$
    \Send{$\msg{Prepare}(v_r, n_r, (c, s, \mathit{op}), k_r)$}{others}
  \ElsIf{$s = C_r[c].s$ and $C_r[c]$ holds a reply $x$}
    \State reply $(v_r, c, s, x)$
  \EndIf
\On{$\msg{Prepare}(v, n, e, k)$} \why{lem:h-prepare}
  \Guard{\Call{Accept}{$v$}}
  \If{$n > n_r + 1$} \Call{StateTransfer}{}; \Return \EndIf
  \If{$n = n_r + 1$} \Call{Append}{$e$} \EndIf
  \State \Call{CommitUpTo}{$\min(k, n_r)$, false}
  \Send{$\msg{PrepareOk}(v_r, n_r, \mathrm{id}_r)$}{the primary} \why{lem:newcommit}
\On{$\msg{PrepareOk}(v, n, i)$} \why{lem:h-prepareok}
  \Guard{$v = v_r \wedge \mathrm{isPrimary}(r) \wedge \status_r = \normal \wedge n > k_r \wedge n \in \mathrm{dom}\, A_r$}
  \State $A_r[n] \gets A_r[n] \cup \{i\}$
  \If{$i$ was new and $|A_r[n]| = Q$}
    \State \Call{CommitUpTo}{$n$, true};\ remove every $n' \le n$ from $A_r$
  \EndIf
\On{$\msg{Commit}(v, k)$} \why{lem:h-commit}
  \Guard{\Call{Accept}{$v$}}
  \If{$k > n_r$} \Call{StateTransfer}{} \Else\ \Call{CommitUpTo}{$k$, false} \EndIf
\On{$\msg{GetState}(i, v, n)$} \why{lem:h-getstate}
  \Guard{$\status_r = \normal \wedge v = v_r \wedge n \le n_r$}
  \Send{$\msg{NewState}(v, L_r(n{..}n_r], n, n_r, k_r)$}{$i$}
\On{$\msg{NewState}(v, L, a, b, k)$} \why{lem:h-newstate}
  \Guard{$v = v_r \wedge |L| = b - a$}
  \If{$\status_r = \transfer$}
    \If{$a > n_r \vee b \le n_r$} \Return \EndIf
    \State \Call{Append}{$e$} for each $e$ in $L(n_r - a{..}b - a]$, in order
    \State \Call{CommitUpTo}{$k$, false};\ $\status_r \gets \normal$
    \Send{$\msg{PrepareOk}(v_r, n_r, \mathrm{id}_r)$}{the primary} \why{lem:newcommit}
  \ElsIf{$\status_r = \viewchange \wedge \mathit{catching}_r \wedge a = k_r$}
    \State \Call{InstallLog}{$L_r(0{..}a] \cdot L$} \dep{app:transfer}
    \State \Call{CommitUpTo}{$k$, false};\ \Call{EnterNormal}{}
    \Send{$\msg{PrepareOk}(v_r, n_r, \mathrm{id}_r)$}{the primary} \why{lem:newcommit}
  \EndIf
\end{algorithmic}
\end{algorithm}

\subsection{Analysis}

\begin{lemma}[Why views agree]
\label{lem:primary}
Let $p = \primary(v)$. Then:
\begin{enumerate}[nosep,label=(\roman*)]
  \item \emph{One tenure.} $p$ is normal in $v$ during a single interval
    of the execution: once $p$ leaves view $v$, by moving to a later view
    or by crashing, it never returns to $v$.
  \item \emph{Append only.} During that interval $L_p$ only grows, and
    afterwards, until $p$ recovers or moves on, it is frozen.
  \item \emph{Everything comes from it.} Every fragment of $v$ in the
    history is a prefix of the log $p$ held at the time the fragment was
    sent.
\end{enumerate}
Hence all fragments of $v$ are prefixes of one sequence, and any two
agree wherever they overlap, which is \OneLog.
\end{lemma}
\begin{proof}
(i) is clause (a): a replica's view number never decreases, including
across recovery, where it restarts at the persisted view. (ii): while
normal, $p$ changes its log only in \Call{Append}{} on a \msg{Request};
\Call{InstallLog}{} runs only on entering a later view or on recovery,
which is clause (e). (iii): a \msg{Prepare} or \msg{NewState} of $v$ is
sent by $p$ from its own log, or by a backup from a log that \OneLog{}
says agrees with $p$'s; a \msg{StartView} of $v$ is $p$'s log when it
started $v$; a vote from $v$ is the log of a replica last normal in $v$,
which by \OneLog{} agrees with $p$'s.
\end{proof}

\begin{lemma}[A new commit survives]
\label{lem:newcommit}
Suppose $\Inv(\Sigma)$ and replica $r$, normal in $v$ with a log that
agrees with every fragment of $v$, sends $\msg{PrepareOk}(v, n_r, r)$.
If for some $m \le n_r$ this completes a quorum, so that
$\Committed(v, m, L_r[m])$ holds in $\Sigma'$, then in $\Sigma'$ every
fragment of every started view $u > v$ that reaches $m$, and every
non-recovering replica last normal in such a $u$, holds $L_r[m]$ at $m$.
\end{lemma}
\begin{proof}
Strong induction on $u$ over the started views above $v$. Let $u$ be
started from $V$; by the induction hypothesis every vote in $V$ from a
view strictly between $v$ and $u$ holds $L_r[m]$ at $m$, since such a
vote is a fragment of that view. Lemma~\ref{lem:chosen} then gives that
the best vote of $V$ holds $L_r[m]$ at $m$. By clause (b) the
\msg{StartView} of $u$ extends the best vote, so it holds $L_r[m]$ at
$m$; every other fragment of $u$ and every replica last normal in $u$
reaches $m$ by clause (b) and agrees with the \msg{StartView} by
\OneLog.
\end{proof}

\begin{proof}[Proof of Lemma~\ref{lem:h-request}]
The new entry is at op-number $n_r + 1$, and by clause (e) every
fragment of $v_r$ ends at or before $n_r$, so \OneLog{} holds for the new
fragment. The \msg{Prepare} carries $k_r$, backed by \CommitsBacked.
Clause (e) is maintained since the primary's log grew with the fragment.
\end{proof}

\begin{proof}[Proof of Lemma~\ref{lem:h-prepare}]
The appended entry is the fragment $\Frag(v, n-1, [e])$, which agrees
with the replica's log by \OneLog{} for $\ell_r = v_r = v$. The commit is
bounded by the message's $k$, which is backed in $v$ by \CommitsBacked,
and by $n_r$, so $\Backed(\ell_r, k_r)$. The acknowledgement names
$n_r$, which is within the primary's log by clause (e) since every entry
of $L_r$ is a fragment of $v$; clause (f) holds for it since the replica
is normal in $v$ and holds its log. If the acknowledgement completes a
quorum, Lemma~\ref{lem:newcommit} gives \Survival{} for the new
$\Committed$ facts.
\end{proof}

\begin{proof}[Proof of Lemma~\ref{lem:h-prepareok}]
By clause (f) every member of $A_r[n]$ other than the primary sent a
\msg{PrepareOk} for op-number $n$ in $v_r$, and the members are distinct,
so $|A_r[n]| = Q$ gives $\QuorumAcked(v_r, m)$ for all $m \le n$. The
entries at those op-numbers are held by $v_r$ (they are the primary's
own fragments), so $\Committed(v_r, m, L_r[m])$ and $\Backed(v_r, n)$;
\CommitsBacked{} follows for the new $k_r = n$, and the guard of
\Call{CommitUpTo}{} holds since $n \le n_r$ by clause (f).
\end{proof}

\begin{proof}[Proof of Lemma~\ref{lem:h-commit}]
\CommitsBacked{} at the message gives $\Backed(v, k)$, and $v = \ell_r$.
\end{proof}

\begin{proof}[Proof of Lemma~\ref{lem:h-getstate}]
The sender is normal in $v$, so $\ell_r = v$ and by \OneLog{} its log
agrees with every fragment of $v$; the segment sent is a piece of that
log. Its commit number is backed by \CommitsBacked{} at the sender, and
its end $n_r$ reaches the base of $v$ by clause (b).
\end{proof}

\begin{proof}[Proof of Lemma~\ref{lem:h-newstate}]
In \transfer{} status $\ell_r = v_r = v$ by \Local, and the appended
entries are a fragment of $v$ agreeing with its log by \OneLog. While
catching up, the kept prefix $L_r(0{..}k_r]$ is backed in $\ell_r < v$ by
\CommitsBacked, so by \Survival{} the fragment of $v$ holds the same
entries at those op-numbers; and the fragment reaches $k_r$ because by
clause (b) it reaches the base of $v$, which by Lemma~\ref{lem:chosen}
holds every committed entry. The installed suffix is the fragment
itself. Hence the new $L_r$ agrees with every fragment of $v$, which is
\OneLog{} for the new $\ell_r = v$, and $\Backed(v, k_r)$ by
\Survival{} and the message's $k$. The acknowledgement is as in
Lemma~\ref{lem:h-prepare}.
\end{proof}

\section{View change}
\label{sec:viewchange}

\subsection{Overview}

A replica that stops hearing from the primary moves to the next view and
says so. A replica that hears $f$ others say so votes: it sends the new
primary its log tagged with the view the log came from. The new primary
takes the log from the vote with the latest such view, longest on a tie,
which by Lemma~\ref{lem:chosen} holds everything any earlier view
committed, and announces it with \msg{StartView}.
Algorithm~\ref{alg:vchelpers} gives the helpers,
Algorithm~\ref{alg:viewchange} the handlers, and Algorithm~\ref{alg:idle}
the idle step whose timer starts it all.

\begin{algorithm}[tp]
\caption{View-change helpers.}
\label{alg:vchelpers}
\begin{algorithmic}[1]
\Procedure{StartViewChange}{$v$}
  \State $v_r \gets v$;\ $\status_r \gets \viewchange$;\ $\mathit{catching}_r \gets \mathrm{false}$;\ $S_r, V_r \gets \emptyset$;\ $\mathit{voted}_r \gets \mathrm{false}$
  \Send{$\msg{StartViewChange}(v, \mathrm{id}_r)$}{others}
  \State \Call{MaybeVote}{}
\EndProcedure
\Procedure{MaybeVote}{}
  \If{$\status_r = \viewchange \wedge \lnot\mathit{catching}_r \wedge \lnot\mathit{voted}_r \wedge |S_r| \ge f$}
    \State $\mathit{voted}_r \gets \mathrm{true}$;\ \Call{Vote}{}
  \EndIf
\EndProcedure
\Procedure{Vote}{} \why{lem:h-svc}
  \If{$\primary(v_r) = \mathrm{id}_r$} \Call{RecordVote}{$\mathrm{id}_r$, $(\ell_r, L_r, k_r)$}
  \Else\ send $\msg{DoViewChange}(v_r, \mathrm{id}_r, \ell_r, L_r, n_r, k_r)$ to the primary
  \EndIf
\EndProcedure
\Procedure{RecordVote}{$i$, $\mathit{vote}$} \why{lem:h-start}
  \State $V_r[i] \gets \mathit{vote}$
  \If{$|V_r| < Q$} \Return \EndIf
  \State $\mathit{best} \gets$ the vote in $V_r$ with the greatest $(\ell, |L|)$ lexicographically, ties to the highest identity \why{lem:chosen}
  \State $k^* \gets \max\{\, k : (\ell, L, k) \in V_r \,\}$
  \State start $v_r$ from $V_r$ \Comment{adds $(v_r, V_r)$ to $\started$}
  \State \Call{InstallLog}{$\mathit{best}.L$};\ \Call{CommitUpTo}{$k^*$, true};\ \Call{EnterNormal}{}
  \State $A_r \gets \{\, n \mapsto \{\mathrm{id}_r\} : k_r < n \le n_r \,\}$
  \Send{$\msg{StartView}(v_r, L_r, n_r, k_r)$}{others}
\EndProcedure
\end{algorithmic}
\end{algorithm}

\begin{algorithm}[tp]
\caption{View change.}
\label{alg:viewchange}
\begin{algorithmic}[1]
\On{$\msg{StartViewChange}(v, i)$} \why{lem:h-svc}
  \If{$v < v_r$} \Return \EndIf
  \If{$v > v_r$} \Call{StartViewChange}{$v$} \dep{app:viewchange} \EndIf
  \If{$\status_r = \viewchange$} $S_r \gets S_r \cup \{i\}$;\ \Call{MaybeVote}{}
  \ElsIf{$\status_r = \normal \wedge \mathrm{isPrimary}(r)$} send $\msg{StartView}(v_r, L_r, n_r, k_r)$ to $i$
  \EndIf
\On{$\msg{DoViewChange}(v, i, \ell, L, n, k)$} \why{lem:h-svc}
  \Guard{$v \ge v_r \wedge \primary(v) = \mathrm{id}_r \wedge |L| = n$}
  \If{$v > v_r$} \Call{StartViewChange}{$v$} \EndIf
  \If{$\status_r = \normal$} send $\msg{StartView}(v_r, L_r, n_r, k_r)$ to $i$
  \Else\ \Call{RecordVote}{$i$, $(\ell, L, k)$}
  \EndIf
\On{$\msg{StartView}(v, L, n, k)$} \why{lem:h-startview}
  \Guard{$|L| = n \wedge (v > v_r \vee (v = v_r \wedge \status_r = \viewchange))$} \dep{app:viewchange}
  \State $v_r \gets v$;\ \Call{InstallLog}{$L$};\ \Call{CommitUpTo}{$k$, false};\ \Call{EnterNormal}{};\ $A_r \gets \emptyset$
  \Send{$\msg{PrepareOk}(v_r, n_r, \mathrm{id}_r)$}{the primary} \why{lem:newcommit}
\end{algorithmic}
\end{algorithm}

\begin{algorithm}[tp]
\caption{Idle. ``The timer fires'' is decided by the replica's idle
counters, which the environment's schedule drives
(Assumption~\ref{asm:model}).}
\label{alg:idle}
\begin{algorithmic}[1]
\On{idle} \why{lem:h-idle}
  \If{$\status_r = \normal \wedge \mathrm{isPrimary}(r)$}
    \Send{$\msg{Commit}(v_r, k_r)$}{others}
    \State send $\msg{Prepare}(v_r, n, L_r[n], k_r)$ to others for each $n = k_r+1, \dots, n_r$ \dep{app:normal}
  \ElsIf{$\status_r = \normal$}
    \If{the timer fires} \Call{StartViewChange}{$v_r + 1$} \EndIf
  \ElsIf{$\status_r = \recovering$}
    \Send{$\msg{Recovery}(\mathrm{id}_r, x_r, v_r)$}{others}
  \ElsIf{$\status_r = \transfer$}
    \State \Call{StateTransfer}{};\ \textbf{if} the timer fires \textbf{then} \Call{StartViewChange}{$v_r + 1$}
  \ElsIf{the timer fires}
    \State \Call{StartViewChange}{$v_r + 1$}
  \ElsIf{$\mathit{catching}_r$}
    \Send{$\msg{GetState}(\mathrm{id}_r, v_r, k_r)$}{the primary}
  \Else
    \Send{$\msg{StartViewChange}(v_r, \mathrm{id}_r)$}{others};\ \textbf{if} $\mathit{voted}_r$ \textbf{then} \Call{Vote}{}
  \EndIf
\end{algorithmic}
\end{algorithm}

\subsection{Analysis}

\begin{lemma}[The chosen log holds earlier commits]
\label{lem:chosen}
Let $\Inv(\Sigma)$. Suppose $\Holds(v, n, e)$ and a set $X$ of $Q$
distinct identities below $N$ with $\Acked(v, n, x)$ for every $x \in X$.
Let $u > v$, and let $V$ be $Q$ votes for $u$ from distinct replicas,
each satisfying clauses (c) and (d), such that every vote in $V$ from a
view $\ell$ with $v < \ell < u$ holds $e$ at $n$. Then the best vote of
$V$, by $(\ell, |L|)$, holds $e$ at $n$.
\end{lemma}
\begin{proof}
By Lemma~\ref{lem:quorum} some $x \in X$ cast a vote $d = (\ell, L, k)$
in $V$. By clause (d), $\ell \ge v$: $x$ either sent a \msg{PrepareOk} in
$v$ before voting, or is the primary of $v$, which was started since it
holds a fragment. If $\ell = v$: by clause (c) the vote covers $n$, as
an op-number $x$ acknowledged in $v$ or, for the primary, as a fragment
of $v$; and \OneLog(i) gives $L[n] = e$. If $\ell > v$: the hypothesis
gives $L[n] = e$, since $\ell < u$ by clause (a). So $d$ holds $e$ at
$n$. Let $b = (\ell_b, L_b, k_b)$ be the best vote, so
$(\ell_b, |L_b|) \ge (\ell, |L|)$. If $\ell_b = \ell$ then
$|L_b| \ge |L| \ge n$ and $L_b[n] = e$ by \OneLog(i), both being
fragments of $\ell$. If $\ell_b > \ell \ge v$ then $b$ is a vote from a
view strictly between $v$ and $u$, and the hypothesis gives
$L_b[n] = e$.
\end{proof}

\begin{proof}[Proof of Lemma~\ref{lem:h-svc}]
Adopting a larger view preserves clause (a). A vote is
$(\ell_r, L_r, k_r)$ with $\ell_r$ the view $L_r$ came from, so it is
$\Frag(\ell_r, 0, L_r)$, agreeing with $\ell_r$'s other fragments by
\OneLog; its $k_r$ is backed in $\ell_r$ by \CommitsBacked; it is from
$\ell_r < v_r$ by clause (a); by clause (f) it covers every op-number
$r$ acknowledged in $\ell_r$, which is clause (c), and if $r$ is the
primary of $\ell_r$ it covers every fragment of $\ell_r$ by clause (e);
and every acknowledgement $r$ made in a view $v' < v_r$ was made while
$\ell_r \ge v'$, since $\ell_r$ only rises, which is clause (d). A
\msg{StartView} sent by a normal primary in reply is its current log, a
fragment of $v_r$ by clause (e).
\end{proof}

\begin{proof}[Proof of Lemma~\ref{lem:h-start}]
The $Q$ votes are from distinct replicas and $v_r$ was not started
before, as its primary was not normal in it, so clause (b) holds for the
new $(v_r, V)$. For $u < v_r$ and $\Committed(u, n, e)$, every vote in
$V$ from a view between $u$ and $v_r$ holds $e$ at $n$ by \Survival, so
Lemma~\ref{lem:chosen} gives $\mathit{best}.L[n] = e$, which is
\Survival{} for the new \msg{StartView} and for the primary itself. The
installed log is the new view's only fragment, so \OneLog. Each $k$ in
$V$ is backed in its vote's $\ell < v_r$ by \CommitsBacked, and by
\Survival{} the chosen log holds the same entries, so
$\Backed(v_r, k^*)$; in particular $k^* \le |\mathit{best}.L|$ and
$k_r \le |\mathit{best}.L|$, so both guards hold.
\end{proof}

\begin{proof}[Proof of Lemma~\ref{lem:h-startview}]
The installed log is the fragment of $v$ that the \msg{StartView}
carries, which gives \OneLog{} for $\ell_r = v$. The message's $k$ is
backed in $v$ by \CommitsBacked, and the replica's previous $k_r$,
backed in the old $\ell_r < v$, is held by $v$ at the same op-numbers by
\Survival, so the new $k_r = \max$ is backed and the guard of
\Call{InstallLog}{} holds. The guard refuses a \msg{StartView} for the
current view in normal status, which is what clause (f) needs. The
acknowledgement names $|L|$, within the primary's log by clause (b), and
is as in Lemma~\ref{lem:h-prepare}.
\end{proof}

\begin{proof}[Proof of Lemma~\ref{lem:h-idle}]
The primary's re-sent \msg{Prepare} for op-number $n$ carries $L_r[n]$,
a fragment of $v_r$ that the history already holds by clause (e) and
\OneLog; its \msg{Commit} carries $k_r$, backed by \CommitsBacked. A
re-sent vote is the case of Lemma~\ref{lem:h-svc}. Starting view
$v_r + 1$ is the case of Lemma~\ref{lem:h-svc} with $v > v_r$.
\end{proof}

\section{Recovery}
\label{sec:recovery}

\subsection{Overview}

A replica that comes back with nothing but its persisted view number
asks everyone what it missed. It waits for $Q$ answers to its nonce, one
of them from the primary of the latest view among them, and adopts that
primary's log. Until then it does nothing else. Its \msg{Recovery}
carries the persisted view, so replicas behind it move forward first.
Algorithm~\ref{alg:recovery} gives the handlers.

\begin{algorithm}[tp]
\caption{Recovery.}
\label{alg:recovery}
\begin{algorithmic}[1]
\On{$\msg{Recovery}(i, x, v)$} \why{lem:h-recovery}
  \If{$v > v_r \wedge \status_r \ne \recovering$} \Call{StartViewChange}{$v$}; \Return \EndIf
  \Guard{$\status_r = \normal$}
  \State $\mathit{st} \gets (L_r, k_r)$ if $\mathrm{isPrimary}(r)$ else $\bot$
  \Send{$\msg{RecoveryResponse}(v_r, x, \mathrm{id}_r, \mathit{st})$}{$i$}
\On{$\msg{RecoveryResponse}(v, x, i, \mathit{st})$} \why{lem:h-recovery}
  \Guard{$\status_r = \recovering \wedge x = x_r$}
  \State $R_r[i] \gets (v, \mathit{st})$
  \If{$|R_r| < Q$} \Return \EndIf
  \State $v^* \gets \max\{\, v' : (v', \_) \in \mathrm{ran}\, R_r \,\}$
  \Guard{$v^* \ge v_r \wedge R_r[\primary(v^*)] = (v^*, (L, k))$ for some $L, k$}
  \State $R_r \gets \emptyset$;\ $v_r \gets v^*$;\ \Call{InstallLog}{$L$};\ \Call{CommitUpTo}{$k$, false};\ \Call{EnterNormal}{}
\On{recover with persisted view $v$ and fresh nonce $x$} \whydep{lem:h-recover}{app:recovery}
  \State reset to the initial state, then $\status_r \gets \recovering$;\ $v_r, \ell_r \gets v$;\ $x_r \gets x$
  \Send{$\msg{Recovery}(\mathrm{id}_r, x, v)$}{others}
\end{algorithmic}
\end{algorithm}

\subsection{Analysis}

\begin{proof}[Proof of Lemma~\ref{lem:h-recovery}]
A recovery state is the normal primary's log in $v_r$, a fragment of
$v_r$ by clause (e), which also gives that it covers every
acknowledgement in $v_r$; it answers a \msg{Recovery} with the nonce it
carries; so clause (g). Installing it at the recovering replica makes
$L_r$ that fragment, so \OneLog{} for $\ell_r = v^*$, and $k$ is backed
by \CommitsBacked{} at the message; the guard of \Call{InstallLog}{}
holds since the recovering replica has $k_r = 0$. The recovering replica
held nothing before, so nothing is lost, and $v^* \ge v_r$ keeps clause
(a). A \msg{Recovery} carrying $v > v_r$ moves the receiver to $v$,
which is Lemma~\ref{lem:h-svc}.
\end{proof}

\begin{proof}[Proof of Lemma~\ref{lem:h-recover}]
By Assumption~\ref{asm:model}(iii) the persisted view is the replica's
$v_r$ in $\Sigma$, and by clause (a) in $\Sigma$ no message it sent is
ahead of that, so clause (a) holds after the reset. \Local{} holds for a
recovering replica with an empty log. By Assumption~\ref{asm:model}(iv)
no \msg{RecoveryResponse} answers the new nonce, so clause (g) is
vacuous for it. Every other clause about the replica is either vacuous
for \recovering{} status or unchanged.
\end{proof}

\section{Conclusion}
\label{sec:conclusion}

Viewstamped Replication as published in 2012 is not safe as written, in
two places that lose committed data and one that repeats an operation.
This paper gives the corrected protocol as pseudocode and proves it safe
through an inductive invariant over the history of every message ever
sent, checked in Lean~4. The invariant is the contribution: it is the
precise form of the paper's informal correctness argument, and each
correction to the paper corresponds to a clause of it that the published
version of the protocol would make false. What the proof does not cover
is what clients see, which needs a model of clients and their
retransmissions, and liveness, which needs a fairness assumption and a
different kind of argument.

\subsection*{AI disclosure}
This draft was prepared with Claude Code, an AI coding agent, under the
author's direction. The agent read the cited papers and the mechanised
development and drafted the text, the pseudocode of
Sections~\ref{sec:normal} to~\ref{sec:recovery}, and the proof sketches;
the agent also took part in writing the Lean development. What has been
checked: the Lean proof, by the Lean kernel, with the axioms stated in
Appendix~\ref{app:lean}. What has not: no human has yet checked the Lean
development against the paper's definitions and pseudocode, or the hand
proofs against the Lean, line by line. The correspondence in
Appendix~\ref{app:lean} is the agent's, and a simulated peer review has
already found one place where the pseudocode and the model differ. The
author takes responsibility for the whole.

\bibliography{references}

\appendix

\section{Viewstamped Replication}
\label{app:vr}

This appendix presents the protocol of \citet{liskov2012vr} as
published, and, where Sections~\ref{sec:normal} to~\ref{sec:recovery} correct it, says
why. Section numbers in parentheses refer to that paper.

\subsection{Replicas, views, and quorums}
\label{app:replicas}

A cluster has $N = 2f+1$ replicas with identities $0, \dots, N-1$ and
tolerates $f$ crashes; a quorum is $Q = f+1$ (\S2.2). The cluster moves
through numbered \emph{views}. The primary of view $v$ is replica
$\primary(v) = v \bmod N$; there is no election, only a change of view.

Every replica holds a \emph{view number} $v_r$, a \emph{status} that is
normal, view-change, state-transfer, or recovering, a \emph{log} $L_r$ of
operations, a \emph{commit number} $k_r$ giving the length of the
committed prefix of the log, and the \emph{last normal view} $\ell_r$,
the most recent view in which its status was normal (Figure~2 of the
paper). Entries are numbered by \emph{op-number} from $1$. A \emph{client table} records, per client, the number of its
latest request and, once executed, the reply, so that a retransmitted
request is answered from the table rather than run again.

The paper's replica also has a disk it never writes to. We add one
write: $v_r$ is persisted after every step, before anything the step
produced is sent, and survives a crash. The reason is in
Section~\ref{app:recovery}.

\subsection{Normal operation}
\label{app:normal}

A client sends \msg{Request} to the primary (\S4.1). The primary appends
the operation to its log, assigns it the next op number $n$, records its
own acknowledgement, and sends \msg{Prepare}$(v, n, e, k)$ to the backups,
carrying the entry $e$ and the primary's commit number $k$. A backup
accepts \msg{Prepare} only in view $v_r$, only in normal status, and only
in op-number order; it appends, commits up to $\min(k, |L_r|)$, and
replies \msg{PrepareOk}$(v, |L_r|, r)$, which acknowledges every op up to
its log length. When the primary has $Q$ acknowledgements for op $n$,
counting its own, it commits every op up to $n$ in order, applying each to
the state machine and replying to its client. When idle it sends
\msg{Commit}$(v, k)$ so that backups learn of commits without a new
request.

\begin{departure}{a backup that adopts a lower commit number executes operations twice}
The paper's backups
set their commit number from whatever message arrives. But a
\msg{StartView} can legitimately carry a commit number lower than the
receiver's, since the new primary's log is chosen by last normal view and
length and its commit number is the maximum among the \msg{DoViewChange}
messages it chose from, which may be
behind a backup that committed more in the old view. A backup that
adopts the lower number re-executes the operations between on the next
\msg{Prepare}. We make every commit $k_r := \max(k_r, k)$: the commit
number never decreases, from any message.
\end{departure}

\begin{departure}{a lost message is never re-sent, so the cluster stalls}
The paper ignores re-sending. The
primary re-sends every uncommitted \msg{Prepare} on each idle period
along with \msg{Commit}, and a replica waiting for \msg{NewState} or for
a view change re-sends its request on each idle period. Without this a
single lost message stalls the cluster.
\end{departure}

\subsection{View change}
\label{app:viewchange}

A backup that does not hear from the primary within a timeout advances to
view $v+1$, sets its status to view-change, and sends
\msg{StartViewChange}$(v+1, r)$ to every other replica (\S4.2). A replica
that receives \msg{StartViewChange} for a view above its own moves to
that view and does the same. On $f$ such messages from other replicas for
its current view, a replica sends
\msg{DoViewChange}$(v, r, \ell_r, L_r, |L_r|, k_r)$ to $\primary(v)$,
carrying its log and the view that log is from. On $Q$ of them the new
primary takes the log of the one with the greatest $(\ell, |L|)$ in
lexicographic order, sets its commit number to the greatest $k$ among
them, commits up to it, becomes normal in $v$
with $\ell = v$, and sends \msg{StartView}$(v, L, |L|, k)$ to the others.
A backup that receives \msg{StartView} installs the log, commits up to
$k$ subject to the rule above, becomes normal in $v$, and acknowledges
its whole log with a \msg{PrepareOk}.

The paper's argument (\S8.1) is that a committed op is in $Q$ logs, a
view change reads $Q$ logs, and the two quorums intersect; and that a
replica stops accepting \msg{Prepare} from the old view the moment it
enters view-change status, so the old primary cannot commit anything the
new primary does not see.

\begin{departure}{read as an increment, ``advances its view-number'' lets a primary count votes cast for different views together}
\S4.2 says a replica that
receives \msg{StartViewChange} for a later view ``advances its
view-number'', which can be read as adopting the received view or as
incrementing by one. Under the increment reading a primary can hold
\msg{DoViewChange} messages for views $v$, $v+1$, and $v+2$ at once and
count them together, which builds a log from messages that were not all
for the same view. We
adopt the received view number and increment only on a timeout.
\end{departure}

\begin{departure}{a StartView for the current view would discard entries the replica has acknowledged}
A replica in
normal status in view $v$ must not accept a \msg{StartView} for $v$. Such
a message is one it has already acted on, or a duplicate, and installing
its log would discard entries the replica has since accepted and
acknowledged. We accept a \msg{StartView} for the current view only in
view-change status, and for a later view always.
\end{departure}

\begin{departure}{a view-change backoff reset on entering the view lets the previous primary start the next view change}
A view change that
fails, because the new primary is also down, is followed by another with
a longer timeout. We double the timeout per failed attempt. The backoff
must not be reset on entering normal status: if it is, the previous
primary, freshly normal, times out first in the next view and starts a
change that the others, still backing off, then join, round the ring
without end. It is forgotten only after the new view has been stable for
a full primary timeout.
\end{departure}

\subsection{State transfer}
\label{app:transfer}

A replica that receives a \msg{Prepare} past the end of its log, or a
\msg{Commit} past its log, has missed messages. It sends
\msg{GetState}$(r, v, |L_r|)$ to the primary and receives
\msg{NewState}$(v, L, a, b, k)$ with the entries with op-numbers $a+1$
to $b$, which it appends (\S5.2). A replica that learns of a later view than its own
from a normal-case message is in a different situation: the uncommitted
suffix of its log belongs to a view that has ended and may have been
superseded. The paper says it should ``set its op-number to its
commit-number and remove all entries after this from its log'' before
asking for state.

\begin{departure}{truncating the log before state transfer discards an entry the replica has acknowledged}
Suppose backup $r$ has acknowledged
op $n$ in view $v$, so that op $n$ may be committed on the strength of
$r$'s acknowledgement. Before $r$ learns of the commit it hears of view
$v+1$, truncates to its commit number, and loses op $n$. The change to
$v+1$ then fails and a change to $v+2$ follows, in which $r$ sends a
\msg{DoViewChange} with a log that no longer has op $n$. The quorum that
acknowledged op $n$ and the quorum of \msg{DoViewChange} senders in $v+2$
may intersect only in $r$, and then op $n$, which was committed, is gone.
The argument of \S8.1 needs every replica's \msg{DoViewChange} to cover
every op it has acknowledged, and truncation breaks
exactly that. We have a replica catching up with a later view keep its
log, enter view-change status marked as catching up, ask for the log
after its \emph{commit number}, and on \msg{NewState} replace only the
uncommitted suffix. If the new view's log has op $n$, $r$ keeps it; if
not, the new view was chosen from a quorum that did not need $r$'s
acknowledgement, and op $n$ was never committed.
\end{departure}

\subsection{Recovery}
\label{app:recovery}

A replica that restarts with no memory sets its status to recovering and
sends \msg{Recovery}$(r, x, v)$ with a fresh nonce $x$ (\S4.3). Normal
replicas answer \msg{RecoveryResponse} with their view number, and the
primary of that view includes its log and commit number. On $Q$
responses for its nonce, one of them from the primary of the latest view
among them, the replica adopts that primary's state, becomes normal, and
rejoins. While recovering it takes part in nothing else.

The paper argues (\S8.2) that this is safe without a disk because of the
\msg{StartViewChange} round: a replica that has sent \msg{DoViewChange}
did so only after $f$ others announced the view, so if it crashes and
recovers, those others are still in the new view or beyond and it will
recover into it. The round is offered as the equivalent of the write of
the view id to stable storage that the original protocol made at the end
of every view change \citep[\S4]{oki1988vr}.

\begin{departure}{diskless recovery lets a replica forget that it started a view change}
The round moves the
problem one message earlier. A replica can forget it sent
\msg{StartViewChange} as easily as it can forget a \msg{DoViewChange}.
With three replicas: the replica that will be the new primary sends
\msg{StartViewChange}, crashes before anyone receives it, and recovers in
the old view, because both other replicas are still there and answer
with it. The delayed \msg{StartViewChange} then arrives at the third
replica, which advances its view and sends its own \msg{StartViewChange}
and, now having $f$ of them, a \msg{DoViewChange} back. The recovered
replica, which has no memory of wanting the new view, receives both,
counts them, and completes the view change. The old primary knows nothing
of this and keeps committing. We persist $v_r$ after every step; a
recovered replica starts in that view, and its
\msg{Recovery} message carries it so that the others move forward on
receipt. This is what the 1988 protocol did with its view id, and what
the 2012 paper gave up to need no disk; it turns out to be the one write
that cannot be given up.
\end{departure}

\subsection{What is left out}

Reconfiguration (\S7) is out of scope; a lost machine is replaced by one
that recovers into the same replica identity. Client recovery (\S4.5) is
the client's problem. The optimisations of \S5 and \S6, checkpoints,
witnesses, batching, and reads at the primary, are not treated; the
recovery protocol sends the whole log.

\section{The mechanised proof}
\label{app:lean}

The proof of Theorem~\ref{thm:log} is checked in Lean~4. This appendix
states the mechanised theorem, its assumptions, and how the paper's
definitions and lemmas map onto it. Names in \texttt{monospace} are Lean
identifiers.

\subsection{The theorem}

\begin{verbatim}
theorem safety (m : Machine Op Output St) (sm : St) (config : Config)
    (htwo : 2 <= config.replicaCount)
    (s : System Op Output St) (h : Reachable m sm config s) :
    NoPanic s /\ PrefixAgreement s /\ Durability s
\end{verbatim}

together with
\texttt{commitBounded\_of\_reachable} and \texttt{sentWF\_of\_reachable},
which give commit boundedness and \WellFormed{} for every reachable
cluster. The three depend on the axioms \texttt{propext},
\texttt{Quot.sound}, and \texttt{Classical.choice} only, with no
\texttt{sorry}. \texttt{NoPanic} is the last item of clause (h): no
internal guard of \Call{InstallLog}{} or \Call{CommitUpTo}{} has failed.

\texttt{Reachable} is Definition~\ref{def:reachable}: the closure of the
initial cluster under the four steps of Definition~\ref{def:step}, where a
\emph{recover} step requires \texttt{NonceFresh}, that no \msg{Recovery}
in \texttt{sent} carries the nonce. This and \texttt{2 <= replicaCount}
are the only hypotheses; they are Assumption~\ref{asm:model}(i) and (iv).
Assumption~\ref{asm:model}(iii) is built into the \emph{recover} step,
which restarts the replica with the view number it held.
Assumption~\ref{asm:model}(ii) is built into the model of time: the
handlers keep the protocol's idle counters, and the environment chooses
which replica idles when, so every schedule is a reachable execution.

The model's logs are indexed from zero. Op-number $n$ of this paper is
list index $n - 1$ there, and a commit number $k$ means indices $0$ to
$k - 1$; the definitions below are stated in those terms.

\subsection{The model and its correspondence to the code}

The model is a twin of an implementation: one Lean definition per
function of the implementation, with the same control flow, sending to
an outbox and appending replies to a list. Where the implementation
checks a guard inside \Call{InstallLog}{} or \Call{CommitUpTo}{}, the
model sets a flag instead of stopping, so every handler is total and
``no guard fails'' is a property to prove, \texttt{NoPanic}. The model
is held to the implementation by differential testing rather than by
translation: a generator produces seeded traces of cluster steps, both
sides replay each trace and print every replica's status, view, commit
number, log, and applied operations, and every message and reply sent,
after every step, and the test fails on the first differing line. Every
executable twin of an invariant clause is evaluated on the same traces,
which is how a candidate clause was tried before it was proved.

What the mechanised proof does not cover is what the paper does not
claim: what clients see, since the model has no clients and requests
arrive from the environment, and liveness.

\subsection{Correspondence}

\begin{center}
\small
\begin{tabular}{@{}lp{9.2cm}@{}}
\toprule
This paper & Lean \\
\midrule
Definitions~\ref{def:config} to~\ref{def:messages} & \texttt{Config}, \texttt{Replica}, \texttt{Message} \\
Definition~\ref{def:cluster}, \ref{def:step} & \texttt{System}, \texttt{Step}, \texttt{System.step} \\
Lemma~\ref{lem:quorum} & \texttt{Config.quorum\_intersect}, \texttt{filter\_mem\_length} \\
Definition~\ref{def:voteof} & \texttt{DvcOf} \\
Definition~\ref{def:frag}, Holds & \texttt{Frag}, \texttt{Holds} \\
Acknowledged, Committed, Backed & \texttt{Acked}, \texttt{QuorumAcked}, \texttt{Committed}, \texttt{Backed}, \texttt{MsgBacked} \\
\Local & \texttt{LocalInv}, lifted as \texttt{AllLocal} \\
\WellFormed & \texttt{WF}, lifted as \texttt{SentWF} \\
\OneLog & \texttt{OneLogPerView} \\
\CommitsBacked & \texttt{CommitsBacked} \\
\Survival & \texttt{Survives} \\
$\Inv$, Definition~\ref{def:inv} & \texttt{Inv} \\
Lemma~\ref{lem:mono} & \texttt{Frag.mono}, \texttt{Holds.mono}, \texttt{Committed.mono}, \texttt{Backed.mono} \\
Lemma~\ref{lem:local} & \texttt{LocalInv.onMessage}, \texttt{.onIdle}, \texttt{.recover}; \texttt{OutboxWF.*} \\
Lemma~\ref{lem:h-request} & \texttt{Inv.onRequest} \\
Lemma~\ref{lem:h-prepare} & \texttt{Inv.onPrepare} \\
Lemma~\ref{lem:h-prepareok} & \texttt{Inv.onPrepareOk} \\
Lemma~\ref{lem:h-commit} & \texttt{Inv.onCommit} \\
Lemma~\ref{lem:h-getstate} & \texttt{Inv.onGetState} \\
Lemma~\ref{lem:h-newstate} & \texttt{Inv.onNewState} \\
Lemma~\ref{lem:h-svc} & \texttt{Inv.onStartViewChange}, \texttt{.onDoViewChange}, \texttt{.startViewChange}, \texttt{.sendDoViewChange} \\
Lemma~\ref{lem:h-start} & \texttt{Inv.recordDoViewChange} \\
Lemma~\ref{lem:h-startview} & \texttt{Inv.onStartView} \\
Lemma~\ref{lem:h-recovery} & \texttt{Inv.onRecovery}, \texttt{Inv.onRecoveryResponse} \\
Lemma~\ref{lem:h-recover} & \texttt{Inv.recover} \\
Lemma~\ref{lem:h-idle} & \texttt{Inv.onIdle} \\
Lemma~\ref{lem:init} & \texttt{Inv.init} \\
Theorem~\ref{thm:inductive} & \texttt{Inv.step}, \texttt{inv\_of\_reachable} \\
Lemma~\ref{lem:committed} & \texttt{Inv.committed\_entry} \\
Lemma~\ref{lem:later} & \texttt{Inv.survives\_holds} \\
Lemma~\ref{lem:acker} & \texttt{Inv.holds\_of\_acked} \\
Lemma~\ref{lem:chosen} & \texttt{Inv.bestHolds} \\
Lemma~\ref{lem:newcommit} & \texttt{Inv.viewHolds\_of\_newAck} \\
Theorem~\ref{thm:log}(i) & \texttt{Inv.prefixAgreement} \\
Theorem~\ref{thm:log}(ii) & \texttt{Inv.durability} \\
Theorem~\ref{thm:log}(iii) & \texttt{commitBounded\_of\_reachable} \\
\bottomrule
\end{tabular}
\end{center}

Lemma~\ref{lem:primary}, why views agree, has no single counterpart: it
is the explanation of \OneLog, whose preservation is one obligation of
each handler lemma above. The clauses of Definition~\ref{def:inv} are
the following fields of \texttt{Inv}.

\begin{center}
\small
\begin{tabular}{@{}p{4.6cm}p{9.6cm}@{}}
\toprule
Clause & Lean fields \\
\midrule
(a) Views only rise & \texttt{MessagesBelowView}, \texttt{DvcBelowView}, \texttt{ViewChangeBehind}, \texttt{DvcBehind}, \texttt{PrimaryStartedNormal} \\
(b) Chosen from a quorum, once & \texttt{StartViewChosen}, \texttt{StartedOnce}, \texttt{ViewExtendsBase} \\
(c) Votes cover acknowledgements & \texttt{DvcCoversAcks}, \texttt{DvcPrimaryCovers} \\
(d) Votes come after acknowledgements & \texttt{DvcAfterAcks}, \texttt{DvcAfterOwnView} \\
(e) Primary is longest & \texttt{PrimaryLongest} \\
(f) Acknowledgements are kept & \texttt{AcksHold}, \texttt{AcksCurrent}, \texttt{AcksStarted} \\
(g) Recovery covers acknowledgements & \texttt{RecoveryCoversAcks}, \texttt{RecoveryResponseNonce}, \texttt{RecoveryStateFromPrimary}, \texttt{RecoveryResponseNotSelf}, \texttt{RecordedRRs} \\
(h) Bookkeeping & \texttt{Ids}, \texttt{ReplicaCount}, \texttt{SenderIds}, \texttt{StartedIds}, \texttt{CatchingUpNotPrimary}, \texttt{TransferNotPrimary}, \texttt{PrimaryToOthers}, \texttt{Covered}, \texttt{ReplicasAgree}, \texttt{StartedViews}, \texttt{StartViewChangePos}, \texttt{RecordedDvcs}, \texttt{Drained}, \texttt{Clean}, \texttt{TwoReplicas}, \texttt{NoPanic} \\
\bottomrule
\end{tabular}
\end{center}

\subsection{How preservation is organised}

A step at one replica produces a new replica, a list of messages, and at
most one started view; \texttt{System.after} is the cluster with those
applied. \texttt{StepOK} collects one obligation per field of
\texttt{Inv}, each about the new replica, the new messages, and the new
started view only; facts about old messages and untouched replicas
transfer by Lemma~\ref{lem:mono}, and \texttt{Inv.after} turns
\texttt{StepOK} into \texttt{Inv} of the new cluster. A handler is a
chain of helpers, and its proof follows the chain through the cluster as
it will be once the replica's outbox is handed over
(\texttt{System.drainOf}), one reusable link per kind of change:
replacing a replica while keeping its log, installing a log, appending an
entry, committing up to a backed bound, sending each kind of message.

The Survival obligation of a new \texttt{Committed} fact is discharged
where the paragraph ``Where a commit is born'' in
Section~\ref{sec:framework} says: in \texttt{Inv.onPrepare},
\texttt{Inv.onNewState}, and \texttt{Inv.onStartView}. The lemma used
there, \texttt{Inv.viewHolds\_of\_newAck}, is a strong induction on the
started view, and each step of it is \texttt{Inv.bestHolds}. The same lemma discharges the
guards of \Call{InstallLog}{} and \Call{CommitUpTo}{} when a view is
started (\texttt{Inv.recordDoViewChange}), when a \msg{StartView} is
adopted (\texttt{Inv.onStartView}), and when a catching-up replica takes
a \msg{NewState} (\texttt{Inv.onNewState}).

Every field of \texttt{Inv} also has an executable twin that is evaluated
on the states of randomly generated executions of the model, and every
clause was tried that way before it was proved. Several of the clauses in
(a) to (h) were found by that check failing, or by a preservation proof
not closing, rather than written down in advance; that is the usual way
an inductive invariant is arrived at, and the reason the paper presents
$\Inv$ as the contribution.

\end{document}
